Dans le repère orthonormal $\left(\mathrm{A};\vectt{AB},\vectt{AD},\vectt{AE}\right)$, on a :
\[
A(0~;~0~;~0),\quad B(1~;~0~;~0),\quad D(0~;~1~;~0),\quad E(0~;~0~;~1).
\]

\begin{enumerate}
	\item Comme I est le milieu de $[EF]$, J le milieu de $[EH]$ et K le milieu de $[AE]$, on obtient :
	\[
	I\left(\dfrac12~;~0~;~1\right),\qquad
	J\left(0~;~\dfrac12~;~1\right),\qquad
	K\left(0~;~0~;~\dfrac12\right).
	\]

	\item \begin{enumerate}
		\item On a :
		les vecteurs $\vectt{AI}$ et $\vectt{AJ}$ ont respectivement pour coordonnées :
		\[
		\begin{pmatrix}\dfrac12\\0\\1\end{pmatrix}
		\quad\text{et}\quad
		\begin{pmatrix}0\\\dfrac12\\1\end{pmatrix}.
		\]
		Donc :
		\[
		\vectt{AI}\cdot\vectt{AJ}
		=\dfrac12\times 0+0\times\dfrac12+1\times 1=1.
		\]

		\item On a :
		\[
		\|\vectt{AI}\|=\sqrt{\left(\dfrac12\right)^2+1^2}=\dfrac{\sqrt5}{2}
		\]
		et, de même,
		\[
		\|\vectt{AJ}\|=\dfrac{\sqrt5}{2}.
		\]
		Ainsi :
		\[
		\cos\left(\widehat{\mathrm{IAJ}}\right)
		=\dfrac{\vectt{AI}\cdot\vectt{AJ}}{\|\vectt{AI}\|\times \|\vectt{AJ}\|}
		=\dfrac{1}{\frac{\sqrt5}{2}\times \frac{\sqrt5}{2}}
		=\dfrac45.
		\]
		Donc :
		\[
		\widehat{\mathrm{IAJ}}\approx 36,9^\circ.
		\]
	\end{enumerate}

	\item \begin{enumerate}
		\item On a :
		le vecteur $\vectt{KC}$ a pour coordonnées :
		\[
		\begin{pmatrix}1\\1\\-\dfrac12\end{pmatrix}.
		\]
		Or :
		\[
		\vectt{KC}\cdot\vectt{AI}
		=1\times\dfrac12+1\times 0-\dfrac12\times 1=0
		\]
		et
		\[
		\vectt{KC}\cdot\vectt{AJ}
		=1\times 0+1\times\dfrac12-\dfrac12\times 1=0.
		\]
		Le vecteur $\vectt{KC}$ est orthogonal à deux vecteurs non colinéaires du plan $(AIJ)$.
		Il est donc normal au plan $(AIJ)$.

		\item Comme le vecteur $\vectt{KC}$ a pour coordonnées $\begin{pmatrix}1\\1\\-\dfrac12\end{pmatrix}$ et est un vecteur normal au plan $(AIJ)$, une équation cartésienne de ce plan est de la forme :
		\[
		x+y-\dfrac12 z+d=0.
		\]
		Le point $A(0~;~0~;~0)$ appartient au plan, donc $d=0$.
		Ainsi, une équation cartésienne du plan $(AIJ)$ est :
		\[
		x+y-\dfrac12 z=0.
		\]
	\end{enumerate}

	\item \begin{enumerate}
		\item Le point L est le projeté orthogonal de C sur le plan $(AIJ)$.
		La droite perpendiculaire au plan passant par C a pour vecteur directeur $\vectt{KC}$.
		Comme $C(1~;~1~;~0)$, une représentation paramétrique de cette droite est :
		\[
		\begin{cases}
		x=1+t\\
		y=1+t\\
		z=-\dfrac12t
		\end{cases}
		\quad t\in\R.
		\]
		Le point L appartient aussi au plan $(AIJ)$, donc :
		\[
		(1+t)+(1+t)-\dfrac12\left(-\dfrac12t\right)=0.
		\]
		Ainsi :
		\[
		2+2t+\dfrac14t=0,
		\]
		d'où
		\[
		2+\dfrac94t=0
		\quad\text{et donc}\quad
		t=-\dfrac89.
		\]
		Par conséquent :
		\[
		L\left(1-\dfrac89~;~1-\dfrac89~;~-\dfrac12\times\left(-\dfrac89\right)\right),
		\]
		soit :
		\[
		L\left(\dfrac19~;~\dfrac19~;~\dfrac49\right).
		\]

		\item La distance du point C au plan $(AIJ)$ est la longueur $CL$.
		On a :
		le vecteur $\vectt{CL}$ a pour coordonnées :
		\[
		\begin{pmatrix}-\dfrac89\\-\dfrac89\\\dfrac49\end{pmatrix}.
		\]
		Donc :
		\[
		CL=\sqrt{\left(-\dfrac89\right)^2+\left(-\dfrac89\right)^2+\left(\dfrac49\right)^2}
		=\sqrt{\dfrac{64}{81}+\dfrac{64}{81}+\dfrac{16}{81}}
		=\sqrt{\dfrac{144}{81}}
		=\dfrac43.
		\]
		La distance du point C au plan $(AIJ)$ est donc égale à $\dfrac43$.
	\end{enumerate}

	\item \begin{enumerate}
		\item On a :
		\[
		I\left(\dfrac12~;~0~;~1\right)
		\quad\text{et}\quad
		M(1~;~m~;~1).
		\]
		Donc :
		le vecteur $\vectt{IM}$ a pour coordonnées :
		\[
		\begin{pmatrix}\dfrac12\\m\\0\end{pmatrix}.
		\]
		Une représentation paramétrique de la droite $(IM)$ est donc :
		\[
		\begin{cases}
		x=\dfrac12+\dfrac12s\\
		y=ms\\
		z=1
		\end{cases}
		\quad s\in\R.
		\]

		\item La droite $(KC)$ a pour représentation paramétrique :
		\[
		\begin{cases}
		x=t\\
		y=t\\
		z=\dfrac12-\dfrac12t
		\end{cases}
		\quad t\in\R.
		\]
		Les droites $(IM)$ et $(KC)$ ne sont pas parallèles, car leurs vecteurs directeurs
		$\begin{pmatrix}\dfrac12\\m\\0\end{pmatrix}$ et $\begin{pmatrix}1\\1\\-\dfrac12\end{pmatrix}$ ne peuvent pas être colinéaires.

		Elles sont donc coplanaires si et seulement si elles sont sécantes.
		Cherchons une éventuelle intersection. Il faut résoudre :
		\[
		\begin{cases}
		\dfrac12+\dfrac12s=t\\
		ms=t\\
		1=\dfrac12-\dfrac12t.
		\end{cases}
		\]
		La troisième équation donne $t=-1$.
		La première donne alors :
		\[
		\dfrac12+\dfrac12s=-1,
		\]
		donc $s=-3$.
		La deuxième équation devient :
		\[
		-3m=-1,
		\]
		donc :
		\[
		m=\dfrac13.
		\]
		Les droites $(IM)$ et $(KC)$ sont coplanaires uniquement lorsque $m=\dfrac13$.
		On ne peut donc pas affirmer qu'elles sont coplanaires quelle que soit la valeur de $m$.
	\end{enumerate}
\end{enumerate}
